\chapter{Metodologia}\label{cap:metodologia}

Este capítulo apresenta o método proposto para comparar uma arquitetura de agente único e uma arquitetura multiagente em atendimentos textuais de um provedor de Internet. A comparação será realizada em condições controladas, utilizando os mesmos cenários conversacionais, a mesma base de informações e o mesmo modelo de linguagem. Além da resposta final, serão observados o desenvolvimento da conversa, as decisões de encaminhamento e os recursos consumidos durante cada execução.

\section{Natureza do trabalho}\label{sec:natureza-trabalho}

O trabalho possui natureza aplicada e experimental. A etapa aplicada compreende o desenvolvimento de dois protótipos de atendimento baseados em modelos de linguagem de grande escala. A etapa experimental consiste na execução de um conjunto comum de cenários conversacionais e na comparação dos resultados obtidos pelas duas arquiteturas.

Não se pretende desenvolver um sistema comercial completo. Os protótipos serão utilizados como instrumentos de experimentação, com as funções necessárias para receber mensagens, manter o histórico da conversa, produzir respostas e registrar informações sobre cada execução.

\section{Delineamento experimental}\label{sec:delineamento}

O experimento terá como fator de comparação a arquitetura utilizada no atendimento. Serão avaliadas uma arquitetura de agente único e uma arquitetura multiagente. As duas soluções receberão as mesmas entradas e serão submetidas às mesmas condições de execução.

Na arquitetura de agente único, um único componente será responsável por interpretar as mensagens, manter o contexto, consultar as regras disponíveis e produzir as respostas. Na arquitetura multiagente, uma etapa de triagem definirá o primeiro destino da conversa. O agente especializado manterá o atendimento enquanto o assunto permanecer em sua área e solicitará nova decisão de encaminhamento quando identificar uma mudança de assunto.

Para reduzir interferências externas, serão mantidos iguais entre os protótipos o modelo de linguagem, os parâmetros de geração, a base geral de informações e as regras de atendimento. Os prompts serão diferentes apenas no que for necessário para representar as responsabilidades de cada arquitetura. A versão do modelo, os parâmetros utilizados e as versões das bibliotecas serão registrados juntamente com os resultados.

\subsection{Variáveis do experimento}\label{subsec:variaveis-experimento}

A variável independente será a arquitetura de atendimento, com dois níveis: agente único e sistema multiagente. As variáveis dependentes serão as métricas de qualidade, encaminhamento e consumo obtidas em cada conversa.

Serão controlados o modelo de linguagem, os parâmetros de geração, a base de informações, os cenários, a ordem das mensagens e o ambiente de execução. Caso sejam realizadas repetições, a quantidade e as condições de repetição serão definidas após o cenário piloto, considerando a variabilidade observada nas respostas.

\section{Etapas do trabalho}\label{sec:etapas-trabalho}

O desenvolvimento será dividido nas seguintes etapas:

\begin{enumerate}
    \item definir uma base comum de informações e regras de atendimento;
    \item especificar as arquiteturas de agente único e multiagente;
    \item elaborar os prompts necessários para cada arquitetura;
    \item implementar os dois protótipos e o registro das execuções;
    \item elaborar e revisar os cenários conversacionais de teste;
    \item realizar um experimento piloto para verificar o funcionamento dos protótipos e das métricas;
    \item executar o conjunto de testes definido após o piloto;
    \item consolidar e comparar os resultados obtidos.
\end{enumerate}

\section{Materiais e ferramentas}\label{sec:materiais-ferramentas}

Os protótipos serão desenvolvidos em Python. Os cenários de teste e os resultados das execuções serão armazenados em arquivos estruturados no formato JSON. O código-fonte, os prompts e os documentos auxiliares serão versionados com Git.

A avaliação automatizada será implementada inicialmente com a biblioteca DeepEval. O acesso ao modelo de linguagem será realizado por meio de API. O modelo, o provedor, os parâmetros de geração e as versões das ferramentas serão definidos e registrados antes da execução do experimento. Planilhas ou scripts em Python poderão ser utilizados para consolidar as medições e produzir as tabelas e os gráficos da análise.

O ambiente Orbsofty poderá ser utilizado para inspeção visual e comparação com um fluxo de atendimento em uso. Entretanto, as medições principais serão realizadas nos protótipos controlados, evitando que regras operacionais externas interfiram na comparação entre as arquiteturas.

\section{Cenários conversacionais}\label{sec:cenarios}

Os testes serão formados por conversas com múltiplos turnos, nas quais as informações serão apresentadas gradualmente. Cada cenário conterá as mensagens do usuário, o assunto esperado em cada etapa, o encaminhamento esperado na arquitetura multiagente, as informações que deverão permanecer no contexto e os critérios necessários para considerar a tarefa concluída.

O cenário piloto começará com uma solicitação comercial e, durante a mesma conversa, passará a tratar de uma falha de conexão. Esse caso permitirá observar se o sistema mantém as informações anteriores, identifica a mudança de assunto e, quando aplicável, transfere o atendimento ao agente adequado. Após o piloto, serão elaborados cenários adicionais para as áreas definidas no escopo do trabalho.

\section{Procedimento experimental}\label{sec:procedimento-experimental}

Cada cenário será executado separadamente nas duas arquiteturas. As mensagens do usuário serão enviadas na ordem prevista no arquivo de teste, e a resposta produzida em cada turno será armazenada antes do envio da mensagem seguinte. Também serão registrados o agente responsável, as decisões de encaminhamento, o tempo de resposta, o consumo de tokens e o número de chamadas ao modelo.

As saídas serão avaliadas em relação aos critérios definidos para o cenário. A execução piloto será utilizada para verificar se os registros produzidos são suficientes, se as métricas escolhidas são aplicáveis e se há necessidade de repetir cada cenário. Alterações realizadas após o piloto serão documentadas antes da execução do conjunto definitivo de testes.

\section{Métricas e critérios de avaliação}\label{sec:metricas}

A comparação considerará métricas relacionadas à qualidade da resposta, ao comportamento da conversa e ao consumo de recursos. Inicialmente, serão observados:

\begin{itemize}
    \item correção e relevância das respostas;
    \item retenção das informações fornecidas nos turnos anteriores;
    \item conclusão da tarefa prevista no cenário;
    \item correspondência entre o encaminhamento realizado e o encaminhamento esperado;
    \item quantidade de transferências entre agentes;
    \item tempo de resposta por turno e tempo total da conversa;
    \item número de chamadas ao modelo e consumo de tokens.
\end{itemize}

As métricas baseadas em avaliação por modelo serão complementadas por verificações determinísticas sempre que houver um resultado objetivo, como o agente de destino, a presença de uma informação obrigatória ou a ocorrência de uma transferência. Os limites numéricos e critérios de sucesso serão definidos após a execução piloto, evitando a adoção de valores sem evidência experimental.

\section{Privacidade e proteção dos dados}\label{sec:privacidade}

Os cenários conversacionais serão produzidos especificamente para o estudo, com base em situações típicas de atendimento de provedores de Internet. Não serão utilizadas conversas, cadastros ou registros de clientes reais.

Os nomes, documentos, endereços, números de telefone, protocolos e demais informações presentes nos testes serão fictícios. A experiência profissional do autor será utilizada somente como referência para a elaboração das situações de atendimento, sem reprodução de casos que permitam identificar clientes, funcionários ou empresas. Dessa forma, não serão enviados dados pessoais reais aos modelos de linguagem ou às ferramentas utilizadas nos experimentos.

\section{Riscos e limitações}\label{sec:riscos}

Entre os riscos previstos estão a variação das respostas do modelo, alterações ou indisponibilidade da API, custo de execução, diferença excessiva entre os prompts e limitações da avaliação automática de respostas abertas. Também existe o risco de uma métrica baseada em modelo atribuir uma pontuação inadequada mesmo quando o comportamento esperado puder ser verificado objetivamente.

Como medidas de controle, todas as execuções serão registradas, as versões utilizadas serão identificadas e os encaminhamentos serão avaliados por regras determinísticas. Uma amostra das respostas poderá ser revisada manualmente. Caso o DeepEval não represente adequadamente alguma dimensão do experimento, a métrica correspondente será substituída ou complementada por um procedimento documentado. Também será realizada uma revisão dos cenários para impedir a inclusão acidental de dados associados a atendimentos reais.